\section{Numerical results}
\label{sec:results}

\subsection{Evaluation protocol}

Truth and surrogate are compared in a float64 rollout harness using the same
integrating-factor Gauss--Legendre method, dealiased Zakharov products, spectral
projection, output times, and initial conditions. Only the DNO evaluation is
replaced in an unguarded comparison; guarded comparisons additionally apply
the conditional spectral viscosity in \cref{sec:soliton-guard} to the
surrogate trajectory. The primary metric is final-time relative surface error
\begin{equation}
  e_\eta(T)=
  \frac{\norm{\eta_\theta(T)-\eta_\star(T)}_2}
       {\norm{\eta_\star(T)}_2}.
  \label{eq:rollout-metric}
\end{equation}
We report its median and 95th percentile together with the numbers of
nonfinite and finite-but-divergent trajectories; a finite trajectory is called
divergent when $e_\eta(T)>1$. Truth trajectories with a nonfinite state or
relative Hamiltonian drift above $10^{-3}$ are excluded before scoring.  We
report separately the number of attempted rollouts and the number with a
valid reference trajectory.

The full C16 evaluation contains 32 initial conditions from each of ten panels
spanning five physical wave families, for 320 trajectories total.  Historical
archive names append \texttt{g0}, \texttt{g1}, or \texttt{modal} to some
panels.  These suffixes mark independent generation shards, not different DNO
orders or physical regimes.  We therefore call them replicates A--C below.  The
evaluation uses no runtime cap, hard
learned-output cutoff, adaptive damping, or case-specific intervention.
Thus its finiteness result does not rely on rollout-time post-processing. The guarded C18
experiment is reported separately because the conditional spectral viscosity
of \cref{sec:soliton-guard} is an intentional numerical intervention. Both
experiments otherwise use the same float64 harness and cached truth.

The linear and two Stokes panels originate from a legacy archive with
$L_{\mathrm{old}}=164$.  Before rollout we apply the fixed-gravity change of
variables derived in \cref{app:rescaling}, with $a=164/(2\pi)$, and then
regenerate the truth trajectory on the canonical $L=2\pi$ grid.  Thus
$h_{\mathrm{old}}=1$ becomes
$h^\sharp\approx0.0383121$ for the linear and finite-depth Stokes panels, while
$h_{\mathrm{old}}=1000$ becomes $h^\sharp\approx38.3121$ for the deep-water
Stokes panel.  No reported rollout directly deploys the learned operator on a
noncanonical period.

\subsection{Baseline accuracy and the rollout discrepancy}

The original CS-DNO, which learned all nonlinear corrections on top of
$\G_0$, reached a best single-step validation loss of
$9.28\times10^{-4}$ on the v3 data. A 13.1-million-parameter FNO reached
$3.69\times10^{-3}$ under a matched loss and data split, despite twice as many
epochs. Thus the Craig--Sulem factorization improved one-step accuracy by a
factor of four with approximately 11.6 times fewer parameters.

This advantage did not eliminate long-time failures.  To distinguish value
error from sensitivity to rollout displacement, write
$z_n^\star=(\eta_n^\star,\xi_n^\star)$ for a truth state,
$z_n=(\eta_n,\xi_n)$ for the corresponding surrogate state, and
$\mathcal Q_\theta(z)=\Gtheta(\eta;h)\xi$.  Let
$\mathcal Q^\star$ denote the order-six reference evaluation used by the
rollout harness.  Define
\begin{equation}
  b_n=\mathcal Q_\theta(z_n^\star)-\mathcal Q^\star(z_n^\star),
  \qquad
  i_n=\mathcal Q_\theta(z_n)-\mathcal Q_\theta(z_n^\star),
  \qquad
  \Gamma_n=\frac{\norm{i_n}_2}{\norm{z_n-z_n^\star}_2}.
  \label{eq:rollout-secant-diagnostic}
\end{equation}
These diagnostics use the nondimensional variables of the evaluation and the
equal-weight product norm
$\norm{(\delta\eta,\delta\xi)}_2
=(\norm{\delta\eta}_2^2+\norm{\delta\xi}_2^2)^{1/2}$.
Thus $\Gamma_n$ is a within-protocol secant diagnostic, not a
coordinate-invariant operator norm.
Frame-aligned archives showed median relative DNO errors on the truth states of
$1.87\times10^{-4}$ and $1.72\times10^{-4}$ for two eventual Tanaka failures.
Near their last finite frames, $\Gamma_n$ reached $8.11$ and $8.42$, while
$\norm{i_n}_2/\norm{b_n}_2$ reached $2.64\times10^4$ and
$2.41\times10^4$.  At the earlier pre-cascade frames, the induced response was
concentrated in $k=32,\ldots,64$ and subsequently spread through
$k=64,\ldots,128$.

The Hadamard defect separated stable and unstable behavior more clearly than
the value error. The order-six reference passed the discrete identity at
roughly $10^{-5}$ or below. Earlier learned models exhibited order-one or
larger relative defects on pre-cascade failing states, even while retaining
small pointwise operator error. A transfer-kernel probe further identified a
spurious same-sign coupling from $k_{\mathrm{in}}=117$ to
$k_{\mathrm{out}}=115$, the type of first-order interaction canceled exactly by
\eqref{eq:g1-kernel}.

\subsection{Ablations and rejected remedies}
\label{sec:rejected-remedies}

Several interventions helped local symptoms without resolving the mechanism.
Fourfold time-step refinement changed failure times only weakly and
nonmonotonically. Two-times dealiasing of the Zakharov right-hand side was
negligible on healthy trajectories but delayed terminal failure by only
approximately $0.05$--$0.08$ time units on representative failing states.
Hard high-band caps reduced some NaNs but produced bounded, inaccurate
trajectories and worsened tail error. Soft high-band penalties and stronger
stage-gain penalties did not remove the failing mode.

\paragraph{Rejected GL2 stage-tangent experiment.}
For completeness, we define the stage-gain penalty used in these earlier
experiments.  Let $\mathcal T_\vartheta$ be one Picard update of the two-stage
system \eqref{eq:gl2-stages}, where $\vartheta=\theta$ uses the learned DNO and
$\vartheta=\star$ uses the order-six reference DNO.  Starting from
$Y^{(0)}=(v_n,v_n)$, we formed the stopped reference predictor
$\overline Y=\operatorname{stopgrad}(\mathcal T_\star(Y^{(0)}))$.
For a small random physical perturbation $p$, inserted identically into the
two stages as $\boldsymbol p=(p,p)$, and finite-secant length $\epsilon$,
define
\begin{equation}
  w_{\vartheta,s}
  =
  \frac{
    \mathcal T_\vartheta(\overline Y+\epsilon p)_s
    -\mathcal T_\vartheta(\overline Y)_s
  }{\epsilon},
  \qquad s\in\{1,2\}.
  \label{eq:stage-secant}
\end{equation}
The perturbations were normalized in the linear Hamiltonian norm
\begin{equation}
  \norm{(\delta\eta,\delta\xi)}_{E,h}^{2}
  =
  g\norm{\delta\eta}_2^2
  +\ip{\delta\xi}{\G_0(h)\delta\xi}.
  \label{eq:linear-energy-norm}
\end{equation}
Let $P_B$ be the smooth Fourier projector onto the tested band.  The
corresponding gain of stage $s$ was
\begin{equation}
  r_{\vartheta,s}
  =
  \frac{\norm{P_Bw_{\vartheta,s}}_{E,h}}
       {\norm{p}_{E,h}+\epsilon_{\mathrm{norm}}},
\end{equation}
and the penalty was the squared hinge
\begin{equation}
  \mathcal L_{\mathrm{stage\text{-}gain}}
  =
  \mathbb E\!\left[
    \frac12\sum_{s=1}^2
    \left[
      r_{\theta,s}-r_{\star,s}
      -\left(\mu r_{\star,s}+\nu\right)
    \right]_+^2
  \right],
  \qquad [a]_+=\max\{a,0\}.
  \label{eq:stage-gain-loss}
\end{equation}
The expectation averages over training states and random low-to-mid and
mid-to-mid perturbations.  Each $p$ has unit energy norm,
$\epsilon$ is log-uniform on $[10^{-6},10^{-3}]$, and
$\epsilon_{\mathrm{norm}}=10^{-12}$.  A companion response-matching term compared
$P_Bw_{\theta,s}$ directly with $P_Bw_{\star,s}$.  The decisive C8 and C10
runs used $\mu=\nu=0$, so any learned gain above the reference gain was
penalized.  This construction did not assert
$r_{\theta,s}<1$ and therefore did not impose contraction of the Picard map;
it only limited excess gain relative to the discrete reference map.  C8 left
the Tanaka nonfinite count at $3/32$ and worsened both Tanaka and
Benjamin--Feir median errors.  Restricting the test to $64\leq|k|\leq128$ in
C10 reduced the count to $2/32$, but retained adverse finite-error tradeoffs.
We therefore did not include this penalty in
\eqref{eq:combined-loss}.

Fine-tuning on
model-harvested pre-cascade states increased failures in both Tanaka and
Benjamin--Feir tests; inspection showed that the harvested pack already
contained the old model's high-frequency contamination.  A direct
residual-only use of the fixed-$g$ variables from
\eqref{eq:zakharov-rescaling} was also tested: it supplied the learned residual
with $\eta/h$, $\xi/h^{3/2}$, $q/\sqrt h$, and $h|k|$, while leaving the exact
physical $\G_0+\G_1$ backbone unchanged.  On the matched one-percent-data
smoke, validation loss worsened from $0.0131081$ to $0.0142273$; on the shallow
Tanaka case-27 probe, relative DNO error worsened from $0.005059$ to $0.010280$
and relative speed bias from $0.0871\%$ to $0.7394\%$.  We therefore rejected
that particular learned-residual parameterization.  This result does not
contradict the exact continuum identity \eqref{eq:dno-rescaling}; it shows that
imposing the scaling on only one component of a mixed physical/learned model
did not improve optimization.  These tests did not support time-step
refinement, output clipping, or contaminated-state fine-tuning as explanations
or remedies for the observed failure.

\subsection{C16 full-suite result}
\label{sec:c16-full-suite}

The clean C16 replay shares the exact $\G_0+\G_1+O(\eta^2)$ decomposition and
self-adjoint structural constraints of \eqref{eq:current-model}, but uses the
earlier 860,928-parameter capacity and a relative $H^1$ plus sparse-Hadamard
objective. It predates the current $L^2$, mode, and localized-tangent objective
\eqref{eq:combined-loss}. The
best validation checkpoint occurs at
cumulative epoch 39, with validation loss
$6.736928\times10^{-3}$; epoch 40 is slightly worse at
$6.737164\times10^{-3}$. The value is not directly comparable with the old v3
loss because the dataset, architecture, normalization path, and auxiliary
objective differ. Long-time rollout is the relevant test.

\begin{table}[t]
  \centering
  \small
  \caption{Clean C16 epoch-39 full-suite results with no rollout-time
  intervention.  Each panel contains 32 attempted initial conditions.
  Nonfinite counts use all attempts; divergences and error quantiles use only
  the truth-valid cases shown in the second column.}
  \label{tab:c16-results}
  \begin{tabular}{lrrrrr}
    \toprule
    Evaluation panel & Truth-valid & Nonfinite & Divergent
      & Median $e_\eta(T)$ & 95th percentile\\
    \midrule
    Tanaka, replicate A          & 32 & 0/32 & 1/32 & 0.08643 & 0.51111\\
    Tanaka, replicate B          & 32 & 0/32 & 1/32 & 0.11541 & 0.66002\\
    Benjamin--Feir, replicate A  & 32 & 0/32 & 0/32 & 0.01208 & 0.20452\\
    Benjamin--Feir, replicate B  & 31 & 0/32 & 0/31 & 0.01773 & 0.06932\\
    Benjamin--Feir, replicate C  & 31 & 0/32 & 0/31 & 0.01409 & 0.15342\\
    Random sea, deep   & 32 & 0/32 & 0/32 & 0.00283 & 0.02066\\
    Random sea, finite & 32 & 0/32 & 0/32 & 0.00542 & 0.04091\\
    Linear             & 20 & 0/32 & 0/20 & 0.02707 & 0.08647\\
    Stokes, deep       & 32 & 0/32 & 0/32 & 0.00044 & 0.00112\\
    Stokes, finite     & 31 & 0/32 & 0/31 & 0.00343 & 0.02257\\
    \midrule
    Total              & 305 & 0/320 & 2/305 & & \\
    \bottomrule
  \end{tabular}
\end{table}

No instance of the previously observed NaN mechanism occurs in the unguarded
320-case C16 suite in \cref{tab:c16-results}. Because C16 changed the analytic
backbone, residual structure, shape supervision, capacity, and solver path as a
bundle, this observation does not isolate one of those ingredients as the
cause. The two remaining threshold
crossings are both finite Tanaka trajectories, but they do not have the same
cause. Tanaka replicate-A case 27 is principally a translation: its final raw error is
$1.167$, the crest is displaced by eight grid points, and translation
alignment reduces the error to $0.0608$. The predicted and true final crest
amplitudes are $0.003660$ and $0.003707$. Its initial translation-speed bias is
approximately $-0.60\%$, which accumulates secularly while changing the
Hamiltonian only weakly.

The historical two-epoch C18 global-tangent fine-tune used
\begin{equation}
  c_b(q)=-\frac{\langle q,\eta_{b,x}\rangle}{\|\eta_{b,x}\|_2^2},
  \qquad
  \mathcal L_{\mathrm{tan,glob}}
  =\frac1{|\mathcal I_{56}|}\sum_{b\in\mathcal I_{56}}
  \left[
  \frac{c_b(q^\theta)-c_b(q^\star)}
       {\max\{|c_b(q^\star)|,10^{-3}\}}
  \right]^2,
  \label{eq:c18-global-tangent}
\end{equation}
where $\mathcal I_{56}$ contains source-5 and source-6 samples. This is a
single rigid-wave speed fit and is distinct from the localized C27 term
\eqref{eq:localized-tangent-loss}. It removes the C16 replicate-A divergence: its
unguarded median and 95th percentile become $0.07963$ and
$0.48205$, respectively, with 0/32 divergent trajectories. Benjamin--Feir
replicate B
is essentially unchanged (median $0.01653$, 95th percentile $0.07002$, 0/32
divergent). Tanaka replicate B, however, retains one divergence and its 95th percentile
worsens from $0.66002$ to $0.93169$. This failure of transfer motivates a
separate analysis rather than a larger weight on the same scalar loss.

\subsection{The Tanaka replicate-B case-14 mechanism}
\label{sec:case14-diagnosis}

Case 14 has depth $h=0.1776006$ and consists of two nearly equal left-going
crests with dimensionless amplitudes $0.12514$ and $0.11280$, separated by
$0.57896=3.26h$. This geometry is rare, but the orbit is not absent from the
data: its initial $\eta$ is bit-identical to a source trajectory in the v9
dataset, and 163 of that trajectory's 200 saved states occur in the row-wise
training split.

The failure is neither a pure phase error nor an isolated hallucination of the
learned DNO\@. At $T=200$, translation alignment leaves relative surface errors
$0.506$ for C16 and $0.587$ for C18. Their maximum slopes are $1.124$ and
$1.532$, compared with $0.0777$ for truth, and learned-Hamiltonian drift peaks
near 22--23\% around $t=176$. Yet teacher-forced DNO errors remain below
$1.05\times10^{-3}$ for C16 and $1.17\times10^{-3}$ for C18 over all 251
truth frames. Evaluating the exact order-six DNO on the \emph{predicted}
states agrees with the learned output to roughly 1--7\% through $t=184$.
Thus the dominant high-frequency DNO response is induced by an already
contaminated state; pointwise accuracy on the truth orbit does not control the
transverse numerical trajectory.

The precursor is visible in \eqref{eq:guard-amplitude}. For case 14,
$A(0)=4.5345\times10^{-9}$ and the truth remains at this scale. In contrast,
the predicted amplitudes evolve as follows.

\begin{table}[t]
  \centering
  \caption{Production-band amplitude $A(t)$ for Tanaka replicate-B case 14.}
  \label{tab:case14-band-growth}
  \begin{tabular}{rrr}
    \toprule
    $t$ & C16 & C18\\
    \midrule
     20 & $2.34\times10^{-7}$ & $1.94\times10^{-7}$\\
     40 & $8.99\times10^{-7}$ & $7.58\times10^{-7}$\\
    120 & $3.39\times10^{-6}$ & $3.04\times10^{-6}$\\
    140 & $1.11\times10^{-5}$ & $1.04\times10^{-5}$\\
    200 & $5.30\times10^{-5}$ & $7.54\times10^{-5}$\\
    \bottomrule
  \end{tabular}
\end{table}

The former absolute trigger $A>10^{-4}$ therefore never fired. The corrected
threshold in \eqref{eq:guard-amplitude} begins its transition near $t=30$,
while the relative surface error is still below $0.01$. The C18 scalar tangent
loss is also poorly matched to this interaction: approximately 41\% of its
selected source-5/6 rows contain mixed-direction multicrest states, and 99.2\%
of the initial squared DNO error for case 14 is orthogonal to the global
translation direction. Its initial scalar speed error is only
$9.6\times10^{-5}$ relative. These observations explain why improving the
global scalar diagnostic does not prevent this transverse cascade.

\subsection{Causal guard test and guarded C18 result}
\label{sec:guard-results}

As a causal control, applying a relative-growth trigger and fixed
$k_{\mathrm{eff}}=64$ to the original C16 Tanaka replicate-B evaluation changes its
divergence count from 1/32 to 0/32, its median from $0.1154$ to $0.1046$, and
its 95th percentile from $0.6600$ to $0.5809$. For case 14 specifically, the
final error decreases from $1.1987$ to $0.09835$, maximum slope decreases from
$1.124$ to $0.0790$, and final learned-Hamiltonian drift decreases from
$7.23\%$ to $1.16\%$. This intervention changes the outcome without changing
the checkpoint, which supports the spectral-cascade diagnosis.

The production experiment combines C18 with the intrinsic selector and
depth-scaled guard in \cref{sec:soliton-guard}. Its matched Tanaka results are
shown in \cref{tab:c18-guard-results}.

\begin{table}[t]
  \centering
  \caption{C18 Tanaka results before and after the soliton-selective spectral
  guard. Each row contains 32 scored initial conditions; all remain finite.}
  \label{tab:c18-guard-results}
  \begin{tabular}{llrrr}
    \toprule
    Evaluation panel & Guard & Divergent & Median $e_\eta(T)$ & 95th percentile\\
    \midrule
    Tanaka, replicate A & off & 0/32 & 0.07963 & 0.48205\\
    Tanaka, replicate A & on  & 0/32 & 0.05065 & 0.48206\\
    Tanaka, replicate B & off & 1/32 & 0.11134 & 0.93169\\
    Tanaka, replicate B & on  & 0/32 & 0.11135 & 0.78499\\
    \bottomrule
  \end{tabular}
\end{table}

For case 14, the production guard reduces final error from $1.2875$ to
$0.1573$ and maximum slope from $1.532$ to $0.0794$, close to the truth value
$0.0777$. This is the first single-checkpoint/integrator configuration with no
nonfinite or divergent trajectory in either 32-case Tanaka panel.

Two limitations prevent a stronger claim. First, guarded Tanaka replicate-B case 24
ends at $e_\eta(T)=0.996$, only narrowly below the declared divergence
threshold; zero counted divergences therefore does not imply a comfortable
tail margin. Second, although \eqref{eq:soliton-selector} makes the guard
exactly inactive for the audited non-Tanaka initial conditions, C18 changes
the model parameters. Benjamin--Feir replicate B has been rerun and remains
stable, but the other seven non-Tanaka panels have not yet been evaluated with C18.
Accordingly, the demonstrated claim is 0/64 divergent Tanaka rollouts, not
0/320 for C18.

\subsection{Current C27 result and isolated active-mode-loss ablation}
\label{sec:mode-ablation}

The subsequent C27 model retains the architecture
\eqref{eq:current-model}, replaces the base $H^1$ objective by
\eqref{eq:data-loss}, and uses the complete objective
\eqref{eq:combined-loss}. Both C27 and C28 were evaluated with the same enabled
soliton-selective spectral guard. The selector makes it exactly inactive for
all 256 non-Tanaka initial conditions. At the saved output cadence, no C27
Tanaka state reaches the nominal 50\%-activation trigger, although the archive
cannot exclude activation at an unsaved integrator substep; C28 has not received
the same inactivity audit. Finiteness statements in this subsection are
therefore conditional on this common guarded protocol.

The C27 final epoch-40 checkpoint was evaluated on 32
initial conditions from each of the ten panels.  All saved predicted values
of $\eta$, $\xi$, and $\Gtheta(\eta)\xi$ are finite in all 320 attempted
rollouts.  The reference solver is valid in 305 cases.  On those cases, the
raw terminal error has median, 95th percentile, mean, and maximum
$0.002850$, $0.047694$, $0.011408$, and $0.534156$, respectively.  The numbers
of raw errors above $0.25$, $0.50$, $0.75$, and $1$ are $1$, $1$, $0$, and
$0$.

To separate a displacement from a change of shape, define
\begin{equation}
  e_{\mathrm{align}}
  =\min_{s\in\T_L}
  \frac{\norm{\tau_s\eta_\theta(T)-\eta_\star(T)}_2}
       {\norm{\eta_\star(T)}_2},
  \qquad
  e_{\mathrm{trans}}
  =\sqrt{\max\{e_\eta(T)^2-e_{\mathrm{align}}^2,0\}},
  \label{eq:aligned-errors}
\end{equation}
where $\tau_s f(x)=f(x-s)$.  The shift is obtained by exhaustive periodic
cross-correlation followed by continuous Fourier refinement.  The second
quantity is an alignment-removable diagnostic, not a claim that translation
and deformation form globally orthogonal nonlinear coordinates.

C28 is the isolated deletion test for \eqref{eq:mode-loss}.  It is a fresh
40-epoch replica of C27 with the same initialization seed, data, architecture,
optimizer, base $L^2$ term, localized tangent term, Hadamard term, and schedule.
The serialized configurations differ only in
$\lambda_{\mathrm{mode}}:6\mapsto0$.  Both comparisons use the final epoch-40
checkpoint and byte-identical reference trajectories.  The principal results
are shown in \cref{tab:mode-ablation}.

\begin{table}[t]
  \centering
  \small
  \caption{Exact C27/C28 deletion ablation. Pooled rollout rows use the same
  305 truth-valid cases; the Tanaka replicate-B and pooled Benjamin--Feir rows contain
  32 and 94 cases, respectively. Ratios larger than one favor C27.}
  \label{tab:mode-ablation}
  \begin{tabular}{lrrr}
    \toprule
    Quantity & C27: $\lambda_{\mathrm{mode}}=6$
             & C28: $\lambda_{\mathrm{mode}}=0$ & C28/C27\\
    \midrule
    Validation $\mathcal L_2$       & $0.00092031$ & $0.00089679$ & $0.974$\\
    Raw median                       & $0.002850$   & $0.003263$   & $1.145$\\
    Raw 95th percentile              & $0.047694$   & $0.062459$   & $1.310$\\
    Raw mean                         & $0.011408$   & $0.016560$   & $1.452$\\
    Aligned 95th percentile          & $0.026263$   & $0.027802$   & $1.059$\\
    Alignment-removable 95th percentile
                                      & $0.031974$   & $0.050395$   & $1.576$\\
    Tanaka replicate-B raw 95th percentile
                                      & $0.055637$   & $0.229728$   & $4.129$\\
    Benjamin--Feir raw 95th percentile
                                      & $0.041526$   & $0.063113$   & $1.520$\\
    Benjamin--Feir aligned 95th percentile
                                      & $0.031766$   & $0.026753$   & $0.842$\\
    \bottomrule
  \end{tabular}
\end{table}

Removing the active-mode term improves the one-step base validation loss by
$2.6\%$, but worsens the pooled raw 95th percentile by $31.0\%$, the raw mean
by $45.2\%$, and the alignment-removable 95th percentile by $57.6\%$.  The
raw threshold counts above $0.25/0.50/0.75/1$ change from $1/1/0/0$ to
$3/1/1/1$.  The largest regressions are Tanaka replicate-A case 27
($0.5342\mapsto1.1565$) and Tanaka replicate-B cases 24
($0.03268\mapsto0.45041$) and 29 ($0.00636\mapsto0.35173$).  In the pooled
Benjamin--Feir panel, raw 95th-percentile error increases by $52.0\%$ while
translation-aligned error decreases by $15.8\%$. Because population quantiles
need not be attained by the same cases, this sign difference does not show
that any particular regression is a pure displacement.  It shows only that
the raw and alignment-removable population tails worsen while this aligned
Benjamin--Feir quantile improves.

Every C28 trajectory remains finite under the shared guard, so
\eqref{eq:mode-loss} is not required for finiteness in this matched guarded
benchmark. The isolated conclusion
is instead that, under this fixed training recipe, the term materially
improves the long-time raw and alignment-removable tail that the lower one-step
$\mathcal L_2$ value fails to predict. The pattern is consistent with the
phase/translation mechanism and the
phase-space and modal-rate interpretation in \cref{sec:mode-loss}; the deletion
alone does not prove that weak sidebands are the unique mechanism.  It also
compares one matched seed and coefficients 6 versus 0, so it does not establish
that 6 is optimal.

A separate tangent-off training run, C29, also completed 40 epochs with
$\lambda_{\mathrm{tan}}:10\mapsto0$ as the only training-configuration
change.  A post hoc reconstruction from its archived evaluations gives
C29/C27 ratios of $1.526$ for pooled raw 95th-percentile error, $1.381$ for
pooled raw mean error, and $1.329$ for aligned 95th-percentile error.  The
pooled Tanaka and Benjamin--Feir raw 95th-percentile ratios are $2.873$ and
$1.445$, respectively, and threshold counts above
$0.25/0.50/0.75/1$ change from $1/1/0/0$ to $2/1/1/0$.  All C29 predictions
are finite.  This comparison is not an exact deletion ablation: the recorded
time-integrator source hash changed between the two evaluations, and their
reference-$q$ and protocol hashes are not identical.  It is therefore
secondary evidence for retaining the localized tangent term; a causal
estimate requires replaying both checkpoints under one frozen evaluator and
one reference cache.

\subsection{Reproducibility safeguards}

The C16 result is reported from the corrected replay only. An older
epoch-34--40 continuation is excluded because a device-to-device replication
bug restored the correct optimizer state on one GPU and zero or invalid state
on the other. This made the periodic shape loss fire out of phase and yielded
a nonphysical distributed optimizer trajectory. The clean replay verifies
the full replicated state at startup and after each epoch; throughout the
epoch-17--40 replay, the observed and expected shape-loss activation counts
agree exactly. Both C18 fine-tuning epochs likewise record $373/373$
activations. This distinction is important because the corrupted run happened
to produce some superficially favorable metrics, but cannot support a
scientific comparison.
